\chapter{65}



Dem Gespräch folgten dreissig Minuten Ratlosigkeit. Seine Müdigkeit war
wie weggeblasen. Louis zählte die Runden nicht, die er um die Hütte
zurücklegte. Valentinas Informationen hatten nicht die Klärung gebracht,
die er sich erhofft hatte. Alles war noch viel schlimmer. Seine
Verwirrung hatte sich ins Unermessliche aufgebaut. Über Bastiano und den
Unbekannten durfte er nicht nachdenken, nicht jetzt. Es ging um eine
Entscheidung. Eine gewichtige, wegweisende.

«Mein Gott, was sag ich ihr? Was sag ich ihr?»

Louis sprach die Worte vor sich hin. Blieb zwischendurch stehen und
klopfte immer wieder seine Oberschenkel ab.

Dann setzte er sich vor der Hütte auf die Bank. Er schloss die Augen und
versuchte, sich zu beruhigen.

Also, zu den Fakten. Valentina war bereits im Boot. Sie hatte seine
Handynummer. Er vertraute ihr. Mehr als andern. Er hatte kein
vorsätzliches Verbrechen begangen. Das letzte Detail würde er weglassen.
Valentina würde ihn verstehen. Mit ihr blieb er mit Mailand und dem
weiteren Geschehen dort unten verbunden. Bisher war nicht offiziell
bekannt, dass er geflohen war. Geflohen? Ja, das war er, ausser sich
gewesen und abgehauen.

\sectionbreak

Valentinas Rückruf kam pünktlich.

«Nun, Louis, ich bin zuhause. Erzähl.»

Louis stand auf. Während er kleinschrittig vor der Hütte hin und her ging, 
fand er nach und nach die Worte. Er beschrieb die nächtlichen Szenen,
seine Gefühle, seine Verzweiflung und die Flucht. Er sprach von Giorgias
Zustand. Ihrem irrationalen Verhalten, ihren mysteriösen Aussagen und
ihrem plötzlichen Verschwinden, rücklings über den \emph{Roccia.} Er
wiederholte Giorgias Worte, seine Verwirrung darüber und sprach von
seiner Hilflosigkeit, seiner Angst. Wie er durchdrehte, weil er nicht
verstand, was vor sich ging. Die letzten Sekunden liess er weg. Er
erklärte ihr sein Gefühl, schuldig zu sein und seine Angst, gesucht zu
werden. Und den damit verbundenen Handywechsel.

«Ich habe eine Scheissangst, Valentina. Es sieht so verdammt schlecht
aus für mich. Ich hab Giorgia abstürzen sehen, ich war der Einzige bei
ihr. Noch jetzt sehe ich sie dauernd vor mir. Ich konnte nicht mehr
denken,» Louis atmete schwer, «bin panisch abgehauen, wie ein Irrer, ich
war feige. Was gibt’s für die Polizei Naheliegenderes, als mich zu
verdächtigen? Mich, als verlassenen Freund, der sich schwer tut mit der
Trennung? Ich kann nur hoffen, dass sich das Ganze klärt und Geduld
haben, zu warten.»

Seine letzte Aussage war spontan gekommen. Valentinas Informationen zu
Bastiano und diesem Fremden gaben der Situation eine völlig neue
Perspektive. Auch wenn Louis den Zusammenhang der beiden mit der
nächtlichen Szene nicht herstellen konnte.

Louis’ Verkrampfung löste sich. Valentina hatte still zugehört. Nicht
ein einziges Mal hatte sie ihn unterbrochen. Selbst dann nicht, wenn er
Pausen einlegte. Louis wusste, sie war da. Es tat unendlich gut.

Dass er Mailand verlassen hatte und sich auf einer Schweizer Alp befand,
verschwieg er. Ebenso seinen Identitätswechsel.

Valentina räusperte sich mehrmals. Ihre Betroffenheit erreichte ihn über
die Leitung.

«Mhm, Louis, das ist so, so traurig,» ihre Stimme holperte, es hörte
sich an, als wollte sie die Tränen zurückhalten, es aber nicht schaffte,
«ich hab dich lieb.»

Louis staunte. Er hatte eine andere Reaktion erwartet. Ihre Worte taten
gut. Die Tragödie verlor augenblicklich eine Spur ihrer Schwere.

Als wäre ein Schalter in ihm umgelegt worden, richtete er sich kraftvoll
auf.

«Ich habe dich auch sehr gern, Valentina, ich vertraue dir, wie kaum
jemandem, aber weine nicht, ich, nur ich ganz alleine trage die
Verantwortung dafür. Ich will dich nicht damit belasten, es, es tut mir
leid, wenn…»

«Nein, Louis, schon okay, dafür sind Freundschaften da,» sie redete nun
gefasster weiter, «sag mir, wie ich dir helfen kann.»

«Das hast du schon. Du hast mir zugehört. Es hat mir sehr geholfen.
Danke, Valentina.»

Louis dachte nach. Er musste sich versichern, dass sie ihn nicht
verriet.

«Wie ich schon sagte, bitte ich dich, keinem von unserem Gespräch zu
erzählen und meine Nummer niemandem weiterzugeben. Am besten stellst du
dir vor, du hast sie nie bekommen. Meinst du, dass du das schaffst?»

«Natürlich, Louis,» sie artikulierte sich nun auffallend deutlich,
«etwas möchte ich dir noch sagen.»

«Ja?»

«Die Wahrheit setzt sich durch, immer, manchmal dauert’s. Und vielleicht
kannst du dein Leiden abkürzen, indem du dich stellst.»

Sie war wieder \emph{die} Valentina, die er kannte. Er sah sie vor ihrer
Zitatenwand stehen.

«Danke dir. Ich werd’s mir überlegen. Und bitte, ruf mich nicht an. Ich
melde mich bei dir, okay?»

«Okay, Louis, und, und ich glaube an dich, ich bin für dich da, ciao.»

«Danke, ciao. Ciao, Valentina.»
